\subsection{Tipos de problemas}

Los problemas se caracterizan en tipos con base en la estructura de la salida esperada. Un mismo problema de la vida real se puede modelar de varias formas, que pueden corresponder a diferentes tipos. Algunos tipos son:
\begin{itemize}
  \item \concept{Enumeración}: la salida son todas las soluciones válidas.
  \item \concept{Conteo}: la salida es el número de soluciones válidas.
  \item \concept{Optimización}: la salida es la mejor solución válida de acuerdo a algún criterio.
  \item \concept{Búsqueda}: la salida es alguna solución válida para el problema, si existe.
  \item \concept{Decisión}\label{problema_decision}: la salida es sí o no.
\end{itemize}
Esas categorías no son mutuamente excluyentes, frecuentemente son distintas perspectivas del mismo problema. E.g., el problema de ordenamiento se puede abordar bajo cualquiera de las perspectivas anteriores:
\begin{itemize}
  \item Enumeración: dado un arreglo, listar todas las permutaciones ordenadas de sus elementos.
  \item Conteo: dado un arreglo, determinar cuántas permutaciones distintas de sus elementos están ordenadas.
  \item Optimización: dado un arreglo y un conjunto de operaciones permitidas, encontrar una secuencia de operaciones que lo ordene usando el menor número de operaciones posible.
  \item Búsqueda: dado un arreglo, encontrar alguna permutación de sus elementos que esté ordenada.
  \item Decisión: dado un arreglo, ¿está ordenado?
\end{itemize}
Se pueden formular aún más problemas relacionados a ordenar elementos que encajen en cualquiera de esas categorías. 
